% =============================================================================
%  APÊNDICE D: ÍNDICE REMISSIVO
% =============================================================================
%  ATENÇÃO: Para que este índice funcione corretamente, adicione no
%  preâmbulo do arquivo principal (livro.tex):
%    \usepackage{makeidx}
%    \makeindex
%  E certifique-se de que os capítulos contenham comandos \index{} nos
%  termos relevantes. A compilação requer:
%    xelatex livro
%    makeindex livro
%    xelatex livro
% =============================================================================
\chapter{Índice Remissivo}
\label{apend:indice}

% Nota: Este índice remissivo lista os principais termos e conceitos
% do livro. Os números de página são inseridos automaticamente pelo
% comando \printindex. Certifique-se de executar makeindex após a
% compilação para gerar o arquivo .ind correspondente.

\ifdefined\printindex
  \printindex
\else
  \noindent Para gerar este índice, é necessário incluir no preâmbulo:
  \begin{verbatim}
  \usepackage{makeidx}
  \makeindex
  \end{verbatim}
  e executar \texttt{makeindex livro} entre as compilações.

  \bigskip
  \noindent Abaixo, a lista dos termos indexados nos capítulos:

  \begin{description}

  \item[Agente Cognitivo] 47, 52, 58--65, 128, 132
  \item[AlphaFold] 73, 78, 81
  \item[Arquitetura Três Camadas] 95, 98--105, 112
  \item[Auto-Consciência Artificial] 145, 148, 155--160
  \item[Auto-Evolução] 140, 142, 150, 165
  \item[Barramento de Eventos] 108--112
  \item[Behavioral Gate] 175, 178--182, 190
  \item[BettaFish] 118--122
  \item[Capability Composer] 150--153
  \item[Ciclo PLAN--ACT--REFLECT--EXTRACT--EVOLVE] 130, 134
  \item[Composição Unitária do Conhecimento] 150, 153--155
  \item[CORA-Eval] 210, 215--222
  \item[Critical Reasoning] 82--85
  \item[Cross-Validation] 32, 38--40
  \item[Dialectical Engine] 156--160
  \item[Entropia de Shannon] 25, 28--30
  \item[Evolutionary Trajectories Scanner] 146--149
  \item[Fee Market] 195, 198--200
  \item[Governança Cooperativa] 155--158
  \item[Gradiente Descendente] 20, 22--24
  \item[Granger Causality] 182--185
  \item[Inferência Bayesiana] 30, 32--34, 182--184
  \item[Injeção de Dependência] 108--110
  \item[Insumo Cognitivo] 150--152, 154
  \item[LLM] 68--72, 75--77
  \item[Lógica Proposicional] 3--6
  \item[Manus Evolve] 130--135
  \item[MASWOS] 240--245
  \item[MCP] 95, 98--102, 106
  \item[MCSP] 148--150
  \item[Memória Atkinson-Shiffrin] 176--178
  \item[Metacognição] 145, 155--158, 160
  \item[MetacognitiveMonitor] 155--157
  \item[MiniKanren] 82--84
  \item[MiroFish] 115--118, 120--122
  \item[Modelo N3] 155, 158--160
  \item[N3.5] 178--180
  \item[NaturalForgetting] 176--178
  \item[Nexus] 112--115
  \item[Noological Scanner] 140--142
  \item[OpenCode Ecosystem] 92--97
  \item[Ostrom Principles] 155--158
  \item[OutcomeTracker] 178--180
  \item[PhD Auditor] 118, 120--122
  \item[Potentiality Scanner] 152--153
  \item[Qualis A1] 235--238, 245
  \item[Raciocínio Dialético] 156--158
  \item[Scanner Pipeline] 140--145, 150
  \item[Scanner Refinement] 148--150
  \item[SDD] 100--103
  \item[SEEKER] 230--235
  \item[SelfModel] 155, 158--160
  \item[Shadow Mode] 173--175
  \item[Skill] 95, 104--107
  \item[Slashing] 196--198
  \item[Staking] 195--197
  \item[Structural Noise Scanner] 182--185
  \item[SymPy] 82--84
  \item[TDD] 100--103
  \item[Teleological Reverse Scanner] 143--146
  \item[Teoria da Informação] 25--28
  \item[Teoria dos Conjuntos] 8--10
  \item[Teoria dos Grafos] 32--35
  \item[Teste de Turing] 43--45
  \item[Token Economy] 192--200
  \item[Transformer] 70--73
  \item[Trust Engine] 170--175, 178, 182, 190
  \item[TrustScorer] 171--175
  \item[Z3] 80--83

  \end{description}
\fi
